\documentclass[../main.tex]{subfiles}
\begin{document}

\chapter{Processes and Process Management}
\lessoninfo{
  video={P2L1},
  textbook={OSTEP \cite{ostep} ch.~4--7, 13; Silberschatz et al.\ \cite{silberschatz2018} ch.~3},
  keywords={process, address space, virtual address, page table, process control block, context switch, hot cache, process state, fork, exec, CPU scheduler, timeslice, inter-process communication, shared memory},
}

Lesson~1 described the operating system as the layer that abstracts and
arbitrates the hardware on behalf of applications.  Its most fundamental
abstraction is the process, the representation of an application while it
executes.  This lesson describes what a process is and what state it
comprises, how the operating system represents that state, and how it
manages many processes at once: it switches the CPU between them, creates
them, schedules them, and lets them communicate with one another.

\section{Programs and processes}
\label{sec:l02-process}

An application stored on disk or in flash memory is a
\emph{program}: a static entity, a file of instructions and initial data
that does nothing by itself.\source{P2L1, what is a process; OSTEP
\cite{ostep} ch.~4; Silberschatz et al.\ \cite{silberschatz2018} ch.~3.}
When the program is launched, the operating system loads it into memory and
starts executing it, and it becomes an active entity.

\begin{definition}[Process]
  \label{def:l02-process}
  A \term{process}[プロセス] is an instance of an executing program: the
  state of a program while it executes, loaded in memory.  A process is also
  called a \emph{task}.
\end{definition}

Launching the same program twice creates two processes.  They execute the
same code but each has its own state, for instance different input data or
a different position in the program.  A process represents the execution
state of an \emph{active} application; this does not mean that the process
is using the CPU at every moment.  It may be waiting for user input, or
another process may be running on the CPU while it waits for its turn.

\section{The address space of a process}
\label{sec:l02-address-space}

A process encapsulates all the state of the running application: its code,
its data, and all the variables that the application creates as it
executes.\source{P2L1, what does a process look like, process address
space; OSTEP ch.~4, 13.}  Every element of this state must be uniquely
identified by its address.

\begin{definition}[Address space]
  The \term{address space}[アドレス空間] of a process is the abstraction
  the operating system uses to encapsulate the state of the process: a range
  of addresses from $V_0$ to $V_{\max}$, together with the state stored at
  those addresses.  It is the in-memory representation of the process.
\end{definition}

The state in the address space is of several types, laid out as in
\cref{fig:l02-address-space}:
\begin{description}
  \item[Text and data] The \term{text segment}[テキストセグメント] holds the
    code of the program, the instructions produced by the compiler.  The
    \term{data segment}[データセグメント] holds its static data, such as
    global variables initialized when the program is loaded.  Both are
    static state: they are present from the moment the process is first
    loaded.
  \item[Heap] The \term{heap}[ヒープ] holds state that the process creates
    dynamically during its execution, for example memory allocated with
    \texttt{malloc}.  It grows toward higher addresses.
  \item[Stack] The \term{stack}[スタック] is a dynamic part of the address
    space that grows and shrinks during execution in last-in first-out
    order.  When the process calls a procedure, the state that must be
    restored on return---the return address, saved register values and the
    local variables of the call---is pushed onto the stack; when the
    procedure returns, it is popped.  The stack starts at the high end of
    the address space and grows toward lower addresses.
\end{description}

\begin{figure}[htbp]
  \centering
  % Layout of the address space of a process.
% Usage: % Layout of the address space of a process.
% Usage: % Layout of the address space of a process.
% Usage: \input{figures/l02-address-space}   (width ~112 mm)
\begin{tikzpicture}[x=1mm, y=1mm, font=\footnotesize\sffamily,
  >={Stealth[length=1.8mm]},
  ann/.style={font=\scriptsize\sffamily, anchor=west, align=left, text width=62mm},
  brk/.style={thin, ln-muted}]
  % regions
  \draw[fill=ln-box] (12,0)  rectangle (42,10) node[midway] {\iflangja{テキスト（コード）}{text (code)}};
  \draw[fill=ln-box] (12,10) rectangle (42,20) node[midway] {\iflangja{データ}{data}};
  \draw[fill=ln-box] (12,20) rectangle (42,30) node[midway] {\iflangja{ヒープ}{heap}};
  \draw[fill=white]  (12,30) rectangle (42,50);
  \draw[fill=ln-box] (12,50) rectangle (42,62) node[midway] {\iflangja{スタック}{stack}};
  \node[font=\scriptsize\sffamily, text=ln-muted] at (27,40) {\iflangja{未割り当て}{unallocated}};
  % growth directions
  \draw[->, thick, ln-accent] (19,30) -- (19,36);
  \draw[->, thick, ln-accent] (35,50) -- (35,44);
  % address labels
  \node[anchor=east] at (11,0.8)  {$V_0$};
  \node[anchor=east] at (11,61.2) {$V_{\max}$};
  % annotations
  \draw[brk] (44,0.5) -- (46,0.5) -- (46,19.5) -- (44,19.5);
  \node[ann] at (47,10) {\iflangja{静的な状態：プロセスのロード時から存在する}%
    {static state: present from the moment the process is loaded}};
  \draw[brk] (44,20.5) -- (46,20.5) -- (46,29.5) -- (44,29.5);
  \node[ann] at (47,25) {\iflangja{実行中に動的に確保される；上位アドレスへ伸びる}%
    {allocated dynamically during execution; grows toward higher addresses}};
  \draw[brk] (44,50.5) -- (46,50.5) -- (46,61.5) -- (44,61.5);
  \node[ann] at (47,56) {\iflangja{手続き呼び出しとともに後入れ先出しで伸縮する；下位アドレスへ伸びる}%
    {grows and shrinks last-in first-out with procedure calls; grows toward lower addresses}};
\end{tikzpicture}
   (width ~112 mm)
\begin{tikzpicture}[x=1mm, y=1mm, font=\footnotesize\sffamily,
  >={Stealth[length=1.8mm]},
  ann/.style={font=\scriptsize\sffamily, anchor=west, align=left, text width=62mm},
  brk/.style={thin, ln-muted}]
  % regions
  \draw[fill=ln-box] (12,0)  rectangle (42,10) node[midway] {\iflangja{テキスト（コード）}{text (code)}};
  \draw[fill=ln-box] (12,10) rectangle (42,20) node[midway] {\iflangja{データ}{data}};
  \draw[fill=ln-box] (12,20) rectangle (42,30) node[midway] {\iflangja{ヒープ}{heap}};
  \draw[fill=white]  (12,30) rectangle (42,50);
  \draw[fill=ln-box] (12,50) rectangle (42,62) node[midway] {\iflangja{スタック}{stack}};
  \node[font=\scriptsize\sffamily, text=ln-muted] at (27,40) {\iflangja{未割り当て}{unallocated}};
  % growth directions
  \draw[->, thick, ln-accent] (19,30) -- (19,36);
  \draw[->, thick, ln-accent] (35,50) -- (35,44);
  % address labels
  \node[anchor=east] at (11,0.8)  {$V_0$};
  \node[anchor=east] at (11,61.2) {$V_{\max}$};
  % annotations
  \draw[brk] (44,0.5) -- (46,0.5) -- (46,19.5) -- (44,19.5);
  \node[ann] at (47,10) {\iflangja{静的な状態：プロセスのロード時から存在する}%
    {static state: present from the moment the process is loaded}};
  \draw[brk] (44,20.5) -- (46,20.5) -- (46,29.5) -- (44,29.5);
  \node[ann] at (47,25) {\iflangja{実行中に動的に確保される；上位アドレスへ伸びる}%
    {allocated dynamically during execution; grows toward higher addresses}};
  \draw[brk] (44,50.5) -- (46,50.5) -- (46,61.5) -- (44,61.5);
  \node[ann] at (47,56) {\iflangja{手続き呼び出しとともに後入れ先出しで伸縮する；下位アドレスへ伸びる}%
    {grows and shrinks last-in first-out with procedure calls; grows toward lower addresses}};
\end{tikzpicture}
   (width ~112 mm)
\begin{tikzpicture}[x=1mm, y=1mm, font=\footnotesize\sffamily,
  >={Stealth[length=1.8mm]},
  ann/.style={font=\scriptsize\sffamily, anchor=west, align=left, text width=62mm},
  brk/.style={thin, ln-muted}]
  % regions
  \draw[fill=ln-box] (12,0)  rectangle (42,10) node[midway] {\iflangja{テキスト（コード）}{text (code)}};
  \draw[fill=ln-box] (12,10) rectangle (42,20) node[midway] {\iflangja{データ}{data}};
  \draw[fill=ln-box] (12,20) rectangle (42,30) node[midway] {\iflangja{ヒープ}{heap}};
  \draw[fill=white]  (12,30) rectangle (42,50);
  \draw[fill=ln-box] (12,50) rectangle (42,62) node[midway] {\iflangja{スタック}{stack}};
  \node[font=\scriptsize\sffamily, text=ln-muted] at (27,40) {\iflangja{未割り当て}{unallocated}};
  % growth directions
  \draw[->, thick, ln-accent] (19,30) -- (19,36);
  \draw[->, thick, ln-accent] (35,50) -- (35,44);
  % address labels
  \node[anchor=east] at (11,0.8)  {$V_0$};
  \node[anchor=east] at (11,61.2) {$V_{\max}$};
  % annotations
  \draw[brk] (44,0.5) -- (46,0.5) -- (46,19.5) -- (44,19.5);
  \node[ann] at (47,10) {\iflangja{静的な状態：プロセスのロード時から存在する}%
    {static state: present from the moment the process is loaded}};
  \draw[brk] (44,20.5) -- (46,20.5) -- (46,29.5) -- (44,29.5);
  \node[ann] at (47,25) {\iflangja{実行中に動的に確保される；上位アドレスへ伸びる}%
    {allocated dynamically during execution; grows toward higher addresses}};
  \draw[brk] (44,50.5) -- (46,50.5) -- (46,61.5) -- (44,61.5);
  \node[ann] at (47,56) {\iflangja{手続き呼び出しとともに後入れ先出しで伸縮する；下位アドレスへ伸びる}%
    {grows and shrinks last-in first-out with procedure calls; grows toward lower addresses}};
\end{tikzpicture}

  \caption{Address space of a process.  The heap and the stack grow toward
    each other; the addresses between them are not allocated.}
  \label{fig:l02-address-space}
\end{figure}

\subsection{Virtual and physical addresses}

The addresses $V_0$ to $V_{\max}$ that a process uses need not correspond to
actual locations in the physical memory of the machine.\source{P2L1,
process address space; OSTEP ch.~13.}

\begin{definition}[Virtual and physical address]
  A \term{virtual address}[仮想アドレス] is an address in the address space
  of a process, as used by the process itself.  A
  \term{physical address}[物理アドレス] is the address of a location in the
  physical memory (DRAM) of the machine.
\end{definition}

The memory-management hardware and the operating system translate every
virtual address that the process uses into a physical address.  The mapping
between the two is recorded in a \term{page table}[ページテーブル], which
the operating system maintains for every process.  Decoupling the two kinds
of address frees the layout of the address space from the layout of
physical memory.

\begin{example}
  A process stores a variable \texttt{x} at the virtual address
  \texttt{0x00402A10}.  Its page table maps this address to the physical
  address \texttt{0x1F3A0A10}.  Whenever the process reads or writes
  \texttt{x}, it uses \texttt{0x00402A10}; the translation to
  \texttt{0x1F3A0A10} happens without its involvement.
\end{example}

\subsection{Address space and memory management}

The mapping gives the operating system freedom in how it uses physical
memory, which it needs for two reasons.\source{P2L1, address space and
memory management.}  First, a process rarely uses its whole address space;
large portions, such as the region between the heap and the stack, are
never allocated.  Second, there may not be enough physical memory for all
the state that is allocated.  With 32-bit virtual addresses a single process
can address $2^{32}$ bytes, that is \SI{4}{\gibi\byte}, and several processes
together easily need more memory than the machine has.

The operating system therefore decides dynamically which portions of which
address spaces reside where in physical memory.  Some portions may
temporarily be moved out to disk and brought back into memory when the
process needs them (swapping; Lesson~8).  For every process, the operating
system must know where each part of its address space currently is, in
memory or on disk; this is the information held in the page tables.  The
same information lets it validate memory accesses, that is, check that the
process is actually allowed to access the address it uses.

Every process has its own address space, and the address spaces of
different processes span the same range of virtual addresses.  Two processes
running at the same time can therefore use the very same virtual address,
and their page tables map it to different physical locations
(\cref{fig:l02-address-mapping}).  A virtual address identifies a location
only together with the process that uses it.

\begin{figure}[htbp]
  \centering
  % Two processes with the same virtual address range, mapped by their page
% tables to different physical frames (one region is on disk).
% Usage: % Two processes with the same virtual address range, mapped by their page
% tables to different physical frames (one region is on disk).
% Usage: % Two processes with the same virtual address range, mapped by their page
% tables to different physical frames (one region is on disk).
% Usage: \input{figures/l02-address-mapping}   (width ~112 mm)
\begin{tikzpicture}[x=1mm, y=1mm, font=\footnotesize\sffamily,
  >={Stealth[length=1.6mm]},
  pa/.style={fill=ln-accent!22},  % pages of P1
  pb/.style={fill=ln-source!18},  % pages of P2
  lab/.style={font=\scriptsize\sffamily},
  map/.style={->, thin},
  ttl/.style={font=\footnotesize\sffamily\bfseries, anchor=south}]
  % ---- virtual address spaces (same range for both processes)
  \foreach \x/\s/\n in {5/pa/P1, 91/pb/P2} {
    \draw[\s] (\x,0)  rectangle ++(14,8)  node[midway, lab] {\iflangja{テキスト}{text}};
    \draw[\s] (\x,8)  rectangle ++(14,8)  node[midway, lab] {\iflangja{データ}{data}};
    \draw[\s] (\x,16) rectangle ++(14,8)  node[midway, lab] {\iflangja{ヒープ}{heap}};
    \draw[fill=white] (\x,24) rectangle ++(14,14);
    \draw[\s] (\x,38) rectangle ++(14,10) node[midway, lab] {\iflangja{スタック}{stack}};
    \node[ttl] at (\x+7,49) {\n};
  }
  \node[lab, anchor=east] at (4.5,0.8)  {$V_0$};
  \node[lab, anchor=east] at (4.5,47.2) {$V_{\max}$};
  \node[lab, anchor=west] at (105.5,0.8)  {$V_0$};
  \node[lab, anchor=west] at (105.5,47.2) {$V_{\max}$};
  % ---- page tables
  \draw[fill=white] (25,0) rectangle (35,48)
    node[midway, rotate=90, lab] {\iflangja{P1 のページテーブル}{page table of P1}};
  \draw[fill=white] (75,0) rectangle (85,48)
    node[midway, rotate=90, lab] {\iflangja{P2 のページテーブル}{page table of P2}};
  % ---- physical memory and disk
  \node[ttl] at (55,49) {\iflangja{物理メモリ}{physical memory}};
  \foreach \i/\s in {0/pa, 1/pb, 2/pa, 3/pb, 4/pa, 5/pb, 6/pa} {
    \draw[\s] (47,13+5*\i) rectangle ++(16,5);
  }
  \draw[fill=ln-box, rounded corners=2pt] (47,0) rectangle (63,8)
    node[midway, lab] {\iflangja{ディスク}{disk}};
  % ---- P1 mappings: text->f0, data->f2, heap->f4, stack->f6
  \foreach \y/\f in {4/15.5, 12/25.5, 20/35.5, 43/45.5} {
    \draw (19,\y) -- (25,\y);
    \draw[map] (35,\y) -- (47,\f);
  }
  % ---- P2 mappings: text->disk, data->f1, heap->f3, stack->f5
  \foreach \y/\f in {4/4, 12/20.5, 20/30.5, 43/40.5} {
    \draw (91,\y) -- (85,\y);
    \draw[map] (75,\y) -- (63,\f);
  }
\end{tikzpicture}
   (width ~112 mm)
\begin{tikzpicture}[x=1mm, y=1mm, font=\footnotesize\sffamily,
  >={Stealth[length=1.6mm]},
  pa/.style={fill=ln-accent!22},  % pages of P1
  pb/.style={fill=ln-source!18},  % pages of P2
  lab/.style={font=\scriptsize\sffamily},
  map/.style={->, thin},
  ttl/.style={font=\footnotesize\sffamily\bfseries, anchor=south}]
  % ---- virtual address spaces (same range for both processes)
  \foreach \x/\s/\n in {5/pa/P1, 91/pb/P2} {
    \draw[\s] (\x,0)  rectangle ++(14,8)  node[midway, lab] {\iflangja{テキスト}{text}};
    \draw[\s] (\x,8)  rectangle ++(14,8)  node[midway, lab] {\iflangja{データ}{data}};
    \draw[\s] (\x,16) rectangle ++(14,8)  node[midway, lab] {\iflangja{ヒープ}{heap}};
    \draw[fill=white] (\x,24) rectangle ++(14,14);
    \draw[\s] (\x,38) rectangle ++(14,10) node[midway, lab] {\iflangja{スタック}{stack}};
    \node[ttl] at (\x+7,49) {\n};
  }
  \node[lab, anchor=east] at (4.5,0.8)  {$V_0$};
  \node[lab, anchor=east] at (4.5,47.2) {$V_{\max}$};
  \node[lab, anchor=west] at (105.5,0.8)  {$V_0$};
  \node[lab, anchor=west] at (105.5,47.2) {$V_{\max}$};
  % ---- page tables
  \draw[fill=white] (25,0) rectangle (35,48)
    node[midway, rotate=90, lab] {\iflangja{P1 のページテーブル}{page table of P1}};
  \draw[fill=white] (75,0) rectangle (85,48)
    node[midway, rotate=90, lab] {\iflangja{P2 のページテーブル}{page table of P2}};
  % ---- physical memory and disk
  \node[ttl] at (55,49) {\iflangja{物理メモリ}{physical memory}};
  \foreach \i/\s in {0/pa, 1/pb, 2/pa, 3/pb, 4/pa, 5/pb, 6/pa} {
    \draw[\s] (47,13+5*\i) rectangle ++(16,5);
  }
  \draw[fill=ln-box, rounded corners=2pt] (47,0) rectangle (63,8)
    node[midway, lab] {\iflangja{ディスク}{disk}};
  % ---- P1 mappings: text->f0, data->f2, heap->f4, stack->f6
  \foreach \y/\f in {4/15.5, 12/25.5, 20/35.5, 43/45.5} {
    \draw (19,\y) -- (25,\y);
    \draw[map] (35,\y) -- (47,\f);
  }
  % ---- P2 mappings: text->disk, data->f1, heap->f3, stack->f5
  \foreach \y/\f in {4/4, 12/20.5, 20/30.5, 43/40.5} {
    \draw (91,\y) -- (85,\y);
    \draw[map] (75,\y) -- (63,\f);
  }
\end{tikzpicture}
   (width ~112 mm)
\begin{tikzpicture}[x=1mm, y=1mm, font=\footnotesize\sffamily,
  >={Stealth[length=1.6mm]},
  pa/.style={fill=ln-accent!22},  % pages of P1
  pb/.style={fill=ln-source!18},  % pages of P2
  lab/.style={font=\scriptsize\sffamily},
  map/.style={->, thin},
  ttl/.style={font=\footnotesize\sffamily\bfseries, anchor=south}]
  % ---- virtual address spaces (same range for both processes)
  \foreach \x/\s/\n in {5/pa/P1, 91/pb/P2} {
    \draw[\s] (\x,0)  rectangle ++(14,8)  node[midway, lab] {\iflangja{テキスト}{text}};
    \draw[\s] (\x,8)  rectangle ++(14,8)  node[midway, lab] {\iflangja{データ}{data}};
    \draw[\s] (\x,16) rectangle ++(14,8)  node[midway, lab] {\iflangja{ヒープ}{heap}};
    \draw[fill=white] (\x,24) rectangle ++(14,14);
    \draw[\s] (\x,38) rectangle ++(14,10) node[midway, lab] {\iflangja{スタック}{stack}};
    \node[ttl] at (\x+7,49) {\n};
  }
  \node[lab, anchor=east] at (4.5,0.8)  {$V_0$};
  \node[lab, anchor=east] at (4.5,47.2) {$V_{\max}$};
  \node[lab, anchor=west] at (105.5,0.8)  {$V_0$};
  \node[lab, anchor=west] at (105.5,47.2) {$V_{\max}$};
  % ---- page tables
  \draw[fill=white] (25,0) rectangle (35,48)
    node[midway, rotate=90, lab] {\iflangja{P1 のページテーブル}{page table of P1}};
  \draw[fill=white] (75,0) rectangle (85,48)
    node[midway, rotate=90, lab] {\iflangja{P2 のページテーブル}{page table of P2}};
  % ---- physical memory and disk
  \node[ttl] at (55,49) {\iflangja{物理メモリ}{physical memory}};
  \foreach \i/\s in {0/pa, 1/pb, 2/pa, 3/pb, 4/pa, 5/pb, 6/pa} {
    \draw[\s] (47,13+5*\i) rectangle ++(16,5);
  }
  \draw[fill=ln-box, rounded corners=2pt] (47,0) rectangle (63,8)
    node[midway, lab] {\iflangja{ディスク}{disk}};
  % ---- P1 mappings: text->f0, data->f2, heap->f4, stack->f6
  \foreach \y/\f in {4/15.5, 12/25.5, 20/35.5, 43/45.5} {
    \draw (19,\y) -- (25,\y);
    \draw[map] (35,\y) -- (47,\f);
  }
  % ---- P2 mappings: text->disk, data->f1, heap->f3, stack->f5
  \foreach \y/\f in {4/4, 12/20.5, 20/30.5, 43/40.5} {
    \draw (91,\y) -- (85,\y);
    \draw[map] (75,\y) -- (63,\f);
  }
\end{tikzpicture}

  \caption{Two processes with the same range of virtual addresses.  Their
    page tables map the allocated regions to different physical frames; the
    text of P2 is currently on disk, and the unallocated regions are not
    mapped at all.}
  \label{fig:l02-address-mapping}
\end{figure}

\section{Execution state and the process control block}
\label{sec:l02-pcb}

When the operating system stops a process, it must record what the process
was doing so that it can later resume the process from exactly the same
point.\source{P2L1, process execution state, what is a process control
block; OSTEP ch.~4.}  A program is compiled into a binary, a sequence of
instructions that are not executed strictly in order: jumps, loops and
interrupts change the sequence.  At every moment the CPU must know where in
this sequence the process is.

The \term{program counter}[プログラムカウンタ]\index{PC|see{program counter}}
(PC) holds the address of the next instruction to execute; while the
process runs, it is kept in a register of the CPU.  Other CPU registers hold
data values, addresses, and status information that affects the sequence of
execution.  The \term{stack pointer}[スタックポインタ] (SP) holds the
address of the top of the stack.  Together these values form the execution
state of the process.

\begin{definition}[Process control block]
  A \term{process control block}[プロセス制御ブロック]\index{PCB|see{process control block}}
  (PCB) is a data structure that the operating system maintains for every
  process.  It contains the process state, the process number, the program
  counter and the other CPU registers, the memory limits and the
  virtual-to-physical mappings, the list of open files, the priority, the
  signal mask, and CPU scheduling information.
\end{definition}

The fields of the PCB (\cref{fig:l02-pcb}) are maintained at different
times:
\begin{itemize}
  \item The PCB is created and initialized when the process is created; for
    example, the program counter is set to the first instruction of the
    program.
  \item Some fields are updated when the corresponding state changes.  When
    the process requests more memory, the operating system allocates it and
    updates the memory limits and mappings.
  \item Other fields change far too often to be kept current in the PCB.
    The program counter changes with every instruction, so the CPU maintains
    it in its own register.  The operating system collects the values that
    the CPU keeps for a process and saves them in the PCB whenever the
    process stops running on the CPU.
\end{itemize}

\begin{figure}[htbp]
  \centering
  % Process control blocks of two processes and the CPU registers whose
% contents are saved into and restored from them.
% Usage: % Process control blocks of two processes and the CPU registers whose
% contents are saved into and restored from them.
% Usage: % Process control blocks of two processes and the CPU registers whose
% contents are saved into and restored from them.
% Usage: \input{figures/l02-pcb}   (width ~116 mm)
\begin{tikzpicture}[x=1mm, y=1mm, font=\scriptsize\sffamily,
  >={Stealth[length=1.6mm]},
  fld/.style={draw, minimum width=34mm, minimum height=4.5mm, inner sep=0pt, fill=white},
  cpuf/.style={fld, fill=ln-accent!22},
  reg/.style={draw, minimum width=24mm, minimum height=4.5mm, inner sep=0pt, fill=ln-accent!22},
  ttl/.style={font=\footnotesize\sffamily\bfseries, anchor=south}]
  \foreach \x/\n in {17/1, 99/2} {
    \node[ttl] at (\x,2.5) {\iflangja{P\n の PCB}{PCB of P\n}};
    \node[fld]  at (\x,0)     {\iflangja{プロセス状態}{process state}};
    \node[fld]  at (\x,-4.5)  {\iflangja{プロセス番号}{process number}};
    \node[cpuf] at (\x,-9)    {\iflangja{プログラムカウンタ}{program counter}};
    \node[cpuf] at (\x,-13.5) {\iflangja{レジスタ}{registers}};
    \node[fld]  at (\x,-18)   {\iflangja{メモリの範囲・対応付け}{memory limits / mappings}};
    \node[fld]  at (\x,-22.5) {\iflangja{オープンしたファイルの一覧}{list of open files}};
    \node[fld]  at (\x,-27)   {\iflangja{優先度}{priority}};
    \node[fld]  at (\x,-31.5) {\iflangja{シグナルマスク}{signal mask}};
    \node[fld]  at (\x,-36)   {\iflangja{CPU スケジューリング情報}{CPU scheduling info}};
    \node[fld]  at (\x,-40.5) {\ldots};
  }
  % CPU
  \draw[rounded corners=2pt, fill=ln-box] (46,-19) rectangle (70,2.5);
  \node[font=\footnotesize\sffamily\bfseries] at (58,-2) {CPU};
  \node[reg] at (58,-9)    {PC};
  \node[reg] at (58,-13.5) {\iflangja{レジスタ（SP を含む）}{registers (incl.\ SP)}};
  % save P1, restore P2
  \draw[->, thick] (46,-9) -- (34,-9);
  \draw[->, thick] (46,-13.5) -- (34,-13.5);
  \node[anchor=north, align=center] at (40,-15.5) {\iflangja{P1 の\\状態を\\退避}{save\\state\\of P1}};
  \draw[->, thick] (82,-9) -- (70,-9);
  \draw[->, thick] (82,-13.5) -- (70,-13.5);
  \node[anchor=north, align=center] at (76,-15.5) {\iflangja{P2 の\\状態を\\復元}{restore\\state\\of P2}};
\end{tikzpicture}
   (width ~116 mm)
\begin{tikzpicture}[x=1mm, y=1mm, font=\scriptsize\sffamily,
  >={Stealth[length=1.6mm]},
  fld/.style={draw, minimum width=34mm, minimum height=4.5mm, inner sep=0pt, fill=white},
  cpuf/.style={fld, fill=ln-accent!22},
  reg/.style={draw, minimum width=24mm, minimum height=4.5mm, inner sep=0pt, fill=ln-accent!22},
  ttl/.style={font=\footnotesize\sffamily\bfseries, anchor=south}]
  \foreach \x/\n in {17/1, 99/2} {
    \node[ttl] at (\x,2.5) {\iflangja{P\n の PCB}{PCB of P\n}};
    \node[fld]  at (\x,0)     {\iflangja{プロセス状態}{process state}};
    \node[fld]  at (\x,-4.5)  {\iflangja{プロセス番号}{process number}};
    \node[cpuf] at (\x,-9)    {\iflangja{プログラムカウンタ}{program counter}};
    \node[cpuf] at (\x,-13.5) {\iflangja{レジスタ}{registers}};
    \node[fld]  at (\x,-18)   {\iflangja{メモリの範囲・対応付け}{memory limits / mappings}};
    \node[fld]  at (\x,-22.5) {\iflangja{オープンしたファイルの一覧}{list of open files}};
    \node[fld]  at (\x,-27)   {\iflangja{優先度}{priority}};
    \node[fld]  at (\x,-31.5) {\iflangja{シグナルマスク}{signal mask}};
    \node[fld]  at (\x,-36)   {\iflangja{CPU スケジューリング情報}{CPU scheduling info}};
    \node[fld]  at (\x,-40.5) {\ldots};
  }
  % CPU
  \draw[rounded corners=2pt, fill=ln-box] (46,-19) rectangle (70,2.5);
  \node[font=\footnotesize\sffamily\bfseries] at (58,-2) {CPU};
  \node[reg] at (58,-9)    {PC};
  \node[reg] at (58,-13.5) {\iflangja{レジスタ（SP を含む）}{registers (incl.\ SP)}};
  % save P1, restore P2
  \draw[->, thick] (46,-9) -- (34,-9);
  \draw[->, thick] (46,-13.5) -- (34,-13.5);
  \node[anchor=north, align=center] at (40,-15.5) {\iflangja{P1 の\\状態を\\退避}{save\\state\\of P1}};
  \draw[->, thick] (82,-9) -- (70,-9);
  \draw[->, thick] (82,-13.5) -- (70,-13.5);
  \node[anchor=north, align=center] at (76,-15.5) {\iflangja{P2 の\\状態を\\復元}{restore\\state\\of P2}};
\end{tikzpicture}
   (width ~116 mm)
\begin{tikzpicture}[x=1mm, y=1mm, font=\scriptsize\sffamily,
  >={Stealth[length=1.6mm]},
  fld/.style={draw, minimum width=34mm, minimum height=4.5mm, inner sep=0pt, fill=white},
  cpuf/.style={fld, fill=ln-accent!22},
  reg/.style={draw, minimum width=24mm, minimum height=4.5mm, inner sep=0pt, fill=ln-accent!22},
  ttl/.style={font=\footnotesize\sffamily\bfseries, anchor=south}]
  \foreach \x/\n in {17/1, 99/2} {
    \node[ttl] at (\x,2.5) {\iflangja{P\n の PCB}{PCB of P\n}};
    \node[fld]  at (\x,0)     {\iflangja{プロセス状態}{process state}};
    \node[fld]  at (\x,-4.5)  {\iflangja{プロセス番号}{process number}};
    \node[cpuf] at (\x,-9)    {\iflangja{プログラムカウンタ}{program counter}};
    \node[cpuf] at (\x,-13.5) {\iflangja{レジスタ}{registers}};
    \node[fld]  at (\x,-18)   {\iflangja{メモリの範囲・対応付け}{memory limits / mappings}};
    \node[fld]  at (\x,-22.5) {\iflangja{オープンしたファイルの一覧}{list of open files}};
    \node[fld]  at (\x,-27)   {\iflangja{優先度}{priority}};
    \node[fld]  at (\x,-31.5) {\iflangja{シグナルマスク}{signal mask}};
    \node[fld]  at (\x,-36)   {\iflangja{CPU スケジューリング情報}{CPU scheduling info}};
    \node[fld]  at (\x,-40.5) {\ldots};
  }
  % CPU
  \draw[rounded corners=2pt, fill=ln-box] (46,-19) rectangle (70,2.5);
  \node[font=\footnotesize\sffamily\bfseries] at (58,-2) {CPU};
  \node[reg] at (58,-9)    {PC};
  \node[reg] at (58,-13.5) {\iflangja{レジスタ（SP を含む）}{registers (incl.\ SP)}};
  % save P1, restore P2
  \draw[->, thick] (46,-9) -- (34,-9);
  \draw[->, thick] (46,-13.5) -- (34,-13.5);
  \node[anchor=north, align=center] at (40,-15.5) {\iflangja{P1 の\\状態を\\退避}{save\\state\\of P1}};
  \draw[->, thick] (82,-9) -- (70,-9);
  \draw[->, thick] (82,-13.5) -- (70,-13.5);
  \node[anchor=north, align=center] at (76,-15.5) {\iflangja{P2 の\\状態を\\復元}{restore\\state\\of P2}};
\end{tikzpicture}

  \caption{Process control blocks of P1 and P2.  The shaded fields hold
    values that the CPU maintains in its registers while the process runs;
    they are saved into the PCB when the process stops running and restored
    from it when the process runs again.}
  \label{fig:l02-pcb}
\end{figure}

\section{Context switch}
\label{sec:l02-context-switch}

Consider two processes, P1 and P2, whose PCBs are stored in the memory of
the operating system.\source{P2L1, how is a PCB used, context switch; OSTEP
ch.~6.}  While P1 runs, the CPU registers hold its execution state.  When
the operating system interrupts P1, it saves the contents of the registers
into the PCB of P1 and restores the state of P2 from the PCB of P2 into the
registers; P2 then runs, and P1 is idle.  When the operating system later
interrupts P2, it saves P2's state into its PCB and restores P1's state.
P1 resumes from exactly the point at which it was stopped
(\cref{fig:l02-context-switch}).

\begin{definition}[Context switch]
  A \term{context switch}[コンテキストスイッチ] is the mechanism used by
  the operating system to switch the CPU from the context of one process to
  the context of another.
\end{definition}

\begin{figure}[htbp]
  \centering
  % Context switches between two processes over time.
% Usage: % Context switches between two processes over time.
% Usage: % Context switches between two processes over time.
% Usage: \input{figures/l02-context-switch}   (width ~116 mm)
\begin{tikzpicture}[x=1mm, y=1mm, font=\footnotesize\sffamily,
  >={Stealth[length=1.6mm]},
  ex/.style={draw, fill=ln-accent!22, minimum height=7mm, inner sep=0pt},
  os/.style={draw, fill=ln-box, minimum height=7mm, inner sep=1pt, font=\scriptsize\sffamily, align=center},
  idle/.style={densely dashed, ln-rule, thick},
  lab/.style={font=\scriptsize\sffamily, text=ln-muted}]
  % row labels
  \node[anchor=east] at (11,24) {P1};
  \node[anchor=east] at (11,12) {OS};
  \node[anchor=east] at (11,0)  {P2};
  % P1 executing
  \node[ex, minimum width=24mm] at (24,24) {\iflangja{実行中}{executing}};
  \draw[idle] (36,24) -- (100,24);
  \node[ex, minimum width=16mm] at (108,24) {\iflangja{実行中}{executing}};
  % P2
  \draw[idle] (12,0) -- (56,0);
  \node[ex, minimum width=24mm] at (68,0) {\iflangja{実行中}{executing}};
  \draw[idle] (80,0) -- (116,0);
  % OS work
  \node[os, minimum width=20mm] (s1) at (46,12)
    {\iflangja{PCB\textsubscript{1} へ退避\\PCB\textsubscript{2} から復元}{save into PCB\textsubscript{1}\\restore PCB\textsubscript{2}}};
  \node[os, minimum width=20mm] (s2) at (90,12)
    {\iflangja{PCB\textsubscript{2} へ退避\\PCB\textsubscript{1} から復元}{save into PCB\textsubscript{2}\\restore PCB\textsubscript{1}}};
  % transfers of control
  \draw[->] (36,20.5) -- (36,15.5);
  \draw[->] (56,8.5)  -- (56,3.5);
  \draw[->] (80,3.5)  -- (80,8.5);
  \draw[->] (100,15.5) -- (100,20.5);
  \node[lab, anchor=east] at (35.5,18) {\iflangja{割り込み}{interrupt}};
  \node[lab, anchor=east] at (79.5,6)  {\iflangja{割り込み}{interrupt}};
  % context switch marks and time axis
  \node[lab, anchor=north] at (46,-5) {\iflangja{コンテキストスイッチ}{context switch}};
  \node[lab, anchor=north] at (90,-5) {\iflangja{コンテキストスイッチ}{context switch}};
  \draw[->, ln-muted] (12,-11) -- (116,-11) node[lab, anchor=north east] at (116,-11.5) {\iflangja{時間}{time}};
\end{tikzpicture}
   (width ~116 mm)
\begin{tikzpicture}[x=1mm, y=1mm, font=\footnotesize\sffamily,
  >={Stealth[length=1.6mm]},
  ex/.style={draw, fill=ln-accent!22, minimum height=7mm, inner sep=0pt},
  os/.style={draw, fill=ln-box, minimum height=7mm, inner sep=1pt, font=\scriptsize\sffamily, align=center},
  idle/.style={densely dashed, ln-rule, thick},
  lab/.style={font=\scriptsize\sffamily, text=ln-muted}]
  % row labels
  \node[anchor=east] at (11,24) {P1};
  \node[anchor=east] at (11,12) {OS};
  \node[anchor=east] at (11,0)  {P2};
  % P1 executing
  \node[ex, minimum width=24mm] at (24,24) {\iflangja{実行中}{executing}};
  \draw[idle] (36,24) -- (100,24);
  \node[ex, minimum width=16mm] at (108,24) {\iflangja{実行中}{executing}};
  % P2
  \draw[idle] (12,0) -- (56,0);
  \node[ex, minimum width=24mm] at (68,0) {\iflangja{実行中}{executing}};
  \draw[idle] (80,0) -- (116,0);
  % OS work
  \node[os, minimum width=20mm] (s1) at (46,12)
    {\iflangja{PCB\textsubscript{1} へ退避\\PCB\textsubscript{2} から復元}{save into PCB\textsubscript{1}\\restore PCB\textsubscript{2}}};
  \node[os, minimum width=20mm] (s2) at (90,12)
    {\iflangja{PCB\textsubscript{2} へ退避\\PCB\textsubscript{1} から復元}{save into PCB\textsubscript{2}\\restore PCB\textsubscript{1}}};
  % transfers of control
  \draw[->] (36,20.5) -- (36,15.5);
  \draw[->] (56,8.5)  -- (56,3.5);
  \draw[->] (80,3.5)  -- (80,8.5);
  \draw[->] (100,15.5) -- (100,20.5);
  \node[lab, anchor=east] at (35.5,18) {\iflangja{割り込み}{interrupt}};
  \node[lab, anchor=east] at (79.5,6)  {\iflangja{割り込み}{interrupt}};
  % context switch marks and time axis
  \node[lab, anchor=north] at (46,-5) {\iflangja{コンテキストスイッチ}{context switch}};
  \node[lab, anchor=north] at (90,-5) {\iflangja{コンテキストスイッチ}{context switch}};
  \draw[->, ln-muted] (12,-11) -- (116,-11) node[lab, anchor=north east] at (116,-11.5) {\iflangja{時間}{time}};
\end{tikzpicture}
   (width ~116 mm)
\begin{tikzpicture}[x=1mm, y=1mm, font=\footnotesize\sffamily,
  >={Stealth[length=1.6mm]},
  ex/.style={draw, fill=ln-accent!22, minimum height=7mm, inner sep=0pt},
  os/.style={draw, fill=ln-box, minimum height=7mm, inner sep=1pt, font=\scriptsize\sffamily, align=center},
  idle/.style={densely dashed, ln-rule, thick},
  lab/.style={font=\scriptsize\sffamily, text=ln-muted}]
  % row labels
  \node[anchor=east] at (11,24) {P1};
  \node[anchor=east] at (11,12) {OS};
  \node[anchor=east] at (11,0)  {P2};
  % P1 executing
  \node[ex, minimum width=24mm] at (24,24) {\iflangja{実行中}{executing}};
  \draw[idle] (36,24) -- (100,24);
  \node[ex, minimum width=16mm] at (108,24) {\iflangja{実行中}{executing}};
  % P2
  \draw[idle] (12,0) -- (56,0);
  \node[ex, minimum width=24mm] at (68,0) {\iflangja{実行中}{executing}};
  \draw[idle] (80,0) -- (116,0);
  % OS work
  \node[os, minimum width=20mm] (s1) at (46,12)
    {\iflangja{PCB\textsubscript{1} へ退避\\PCB\textsubscript{2} から復元}{save into PCB\textsubscript{1}\\restore PCB\textsubscript{2}}};
  \node[os, minimum width=20mm] (s2) at (90,12)
    {\iflangja{PCB\textsubscript{2} へ退避\\PCB\textsubscript{1} から復元}{save into PCB\textsubscript{2}\\restore PCB\textsubscript{1}}};
  % transfers of control
  \draw[->] (36,20.5) -- (36,15.5);
  \draw[->] (56,8.5)  -- (56,3.5);
  \draw[->] (80,3.5)  -- (80,8.5);
  \draw[->] (100,15.5) -- (100,20.5);
  \node[lab, anchor=east] at (35.5,18) {\iflangja{割り込み}{interrupt}};
  \node[lab, anchor=east] at (79.5,6)  {\iflangja{割り込み}{interrupt}};
  % context switch marks and time axis
  \node[lab, anchor=north] at (46,-5) {\iflangja{コンテキストスイッチ}{context switch}};
  \node[lab, anchor=north] at (90,-5) {\iflangja{コンテキストスイッチ}{context switch}};
  \draw[->, ln-muted] (12,-11) -- (116,-11) node[lab, anchor=north east] at (116,-11.5) {\iflangja{時間}{time}};
\end{tikzpicture}

  \caption{Two context switches between P1 and P2.  Each process is idle
    (dashed) while the other runs, and its state waits in its PCB.}
  \label{fig:l02-context-switch}
\end{figure}

Context switches are expensive, for two reasons:
\begin{description}
  \item[Direct cost] the CPU cycles spent executing the load and store
    instructions that save the state of one process and restore the state
    of the other.
  \item[Indirect cost] the effect on the processor cache.  A cache access
    takes on the order of a few cycles, whereas an access to main memory
    takes on the order of hundreds of cycles.
\end{description}

While a process runs, much of the data it uses is present in the processor
cache.  Such a cache is a \term{hot cache}[ホットキャッシュ]: most accesses
are served from the cache, and the process performs at its best.  After a
context switch, the next process replaces much of this data with its own.
When the first process resumes, it finds a
\term{cold cache}[コールドキャッシュ]: its data is no longer in the cache,
and its accesses are cache misses that must be served from main memory until
the cache has been refilled.  The operating system should therefore limit
how frequently it switches contexts.

\begin{example}
  Suppose that saving and restoring the registers takes 1\,000 cycles, and
  that an access costs 4 cycles when the data is in the cache and 200 cycles
  when it must be read from main memory.  If P1 accesses 5\,000 data items
  after it resumes that were evicted from the cache while P2 ran, the first
  access to each costs $200 - 4 = 196$ extra cycles, in total
  $5\,000\times 196 = 980\,000$ cycles.  The indirect cost is almost a
  thousand times the direct cost of the switch.
\end{example}

\begin{note}
  A hot cache does not mean that the process may keep the CPU.  A context
  switch is still necessary when, for example, a process of higher priority
  becomes ready or when the running process has used up the time it was
  allotted (\cref{sec:l02-scheduler}).  The cost of the switch is the price
  of sharing the CPU.
\end{note}

\begin{keypoint}
  The PCB holds everything the operating system needs to stop a process and
  resume it later.  A context switch saves the CPU state of one process into
  its PCB and restores another from its PCB; its hidden cost, the cold cache
  it leaves behind, can far exceed the cost of the save and restore.
\end{keypoint}

\section{Process life cycle}
\label{sec:l02-states}

Seen from the CPU, a process alternates between running and not
running.\source{P2L1, process life cycle, process life cycle: states; OSTEP
ch.~4.}  A running process that is interrupted and switched out becomes
idle; it is still ready to execute, and it runs again when the scheduler
dispatches it to the CPU.  A process can be idle for other reasons as well,
and the operating system distinguishes these situations by the
\term{process state}[プロセス状態], which it records in the PCB.  A process
is in one of the following states (\cref{fig:l02-states}):

\begin{description}
  \item[New] A process that is being created is in the
    \term{new state}[新規状態].  The operating system performs
    \emph{admission control}, deciding whether the process can be accepted;
    if so, it allocates and initializes the PCB and some initial
    resources, and the process is \emph{admitted}.
  \item[Ready] A process in the \term{ready state}[実行可能状態] has
    everything it needs to execute except the CPU.  It could run as soon as
    the scheduler selects it; until then it waits for its turn.
  \item[Running] A process in the \term{running state}[実行状態] is
    executing on the CPU; after the \emph{scheduler dispatch} it continues
    from the point stored in its PCB.  A running process leaves this state
    in one of three ways: it is interrupted, for instance because its time
    on the CPU has expired, and returns to the ready state after a context
    switch; it initiates a long operation, such as reading from disk or
    waiting for a keyboard or timer event, and enters the waiting state; or
    it finishes, or encounters an error, and exits.
  \item[Waiting] A process in the \term{waiting state}[待ち状態] cannot
    proceed until some event occurs, such as the completion of an I/O
    operation.  When the event occurs, the process becomes ready again.
  \item[Terminated] A process that has finished all its operations, or has
    encountered an error, calls \emph{exit} and enters the
    \term{terminated state}[終了状態].
\end{description}

\begin{figure}[htbp]
  \centering
  % Process life cycle: states and transitions.
% Usage: % Process life cycle: states and transitions.
% Usage: % Process life cycle: states and transitions.
% Usage: \input{figures/l02-states}   (width ~116 mm)
\begin{tikzpicture}[x=1mm, y=1mm, font=\footnotesize\sffamily,
  st/.style={draw, rounded corners=6pt, fill=ln-box, minimum width=22mm,
             minimum height=8mm, inner sep=2pt},
  >={Stealth[length=1.8mm]},
  lab/.style={font=\scriptsize\sffamily, align=center}]
  \node[st] (new)  at (11,32)  {\iflangja{新規}{new}};
  \node[st] (term) at (105,32) {\iflangja{終了}{terminated}};
  \node[st] (rdy)  at (32,12)  {\iflangja{実行可能}{ready}};
  \node[st, fill=ln-accent!22] (run) at (84,12) {\iflangja{実行}{running}};
  \node[st] (wait) at (58,-12) {\iflangja{待ち}{waiting}};
  \draw[->] (new) -- node[lab, left, pos=0.45] {\iflangja{受け入れ}{admitted}} (rdy);
  \draw[->] (rdy.north east) to[bend left=18]
    node[lab, above] {\iflangja{スケジューラによるディスパッチ}{scheduler dispatch}} (run.north west);
  \draw[->] (run.south west) to[bend left=18]
    node[lab, below] {\iflangja{割り込み}{interrupt}} (rdy.south east);
  \draw[->] (run) -- node[lab, right, pos=0.55] {\iflangja{終了（exit）}{exit}} (term);
  \draw[->] (run) -- node[lab, right, pos=0.6] {\iflangja{I/O・イベント待ち}{I/O or event wait}} (wait);
  \draw[->] (wait) -- node[lab, left, pos=0.4] {\iflangja{I/O・イベント完了}{I/O or event completion}} (rdy);
\end{tikzpicture}
   (width ~116 mm)
\begin{tikzpicture}[x=1mm, y=1mm, font=\footnotesize\sffamily,
  st/.style={draw, rounded corners=6pt, fill=ln-box, minimum width=22mm,
             minimum height=8mm, inner sep=2pt},
  >={Stealth[length=1.8mm]},
  lab/.style={font=\scriptsize\sffamily, align=center}]
  \node[st] (new)  at (11,32)  {\iflangja{新規}{new}};
  \node[st] (term) at (105,32) {\iflangja{終了}{terminated}};
  \node[st] (rdy)  at (32,12)  {\iflangja{実行可能}{ready}};
  \node[st, fill=ln-accent!22] (run) at (84,12) {\iflangja{実行}{running}};
  \node[st] (wait) at (58,-12) {\iflangja{待ち}{waiting}};
  \draw[->] (new) -- node[lab, left, pos=0.45] {\iflangja{受け入れ}{admitted}} (rdy);
  \draw[->] (rdy.north east) to[bend left=18]
    node[lab, above] {\iflangja{スケジューラによるディスパッチ}{scheduler dispatch}} (run.north west);
  \draw[->] (run.south west) to[bend left=18]
    node[lab, below] {\iflangja{割り込み}{interrupt}} (rdy.south east);
  \draw[->] (run) -- node[lab, right, pos=0.55] {\iflangja{終了（exit）}{exit}} (term);
  \draw[->] (run) -- node[lab, right, pos=0.6] {\iflangja{I/O・イベント待ち}{I/O or event wait}} (wait);
  \draw[->] (wait) -- node[lab, left, pos=0.4] {\iflangja{I/O・イベント完了}{I/O or event completion}} (rdy);
\end{tikzpicture}
   (width ~116 mm)
\begin{tikzpicture}[x=1mm, y=1mm, font=\footnotesize\sffamily,
  st/.style={draw, rounded corners=6pt, fill=ln-box, minimum width=22mm,
             minimum height=8mm, inner sep=2pt},
  >={Stealth[length=1.8mm]},
  lab/.style={font=\scriptsize\sffamily, align=center}]
  \node[st] (new)  at (11,32)  {\iflangja{新規}{new}};
  \node[st] (term) at (105,32) {\iflangja{終了}{terminated}};
  \node[st] (rdy)  at (32,12)  {\iflangja{実行可能}{ready}};
  \node[st, fill=ln-accent!22] (run) at (84,12) {\iflangja{実行}{running}};
  \node[st] (wait) at (58,-12) {\iflangja{待ち}{waiting}};
  \draw[->] (new) -- node[lab, left, pos=0.45] {\iflangja{受け入れ}{admitted}} (rdy);
  \draw[->] (rdy.north east) to[bend left=18]
    node[lab, above] {\iflangja{スケジューラによるディスパッチ}{scheduler dispatch}} (run.north west);
  \draw[->] (run.south west) to[bend left=18]
    node[lab, below] {\iflangja{割り込み}{interrupt}} (rdy.south east);
  \draw[->] (run) -- node[lab, right, pos=0.55] {\iflangja{終了（exit）}{exit}} (term);
  \draw[->] (run) -- node[lab, right, pos=0.6] {\iflangja{I/O・イベント待ち}{I/O or event wait}} (wait);
  \draw[->] (wait) -- node[lab, left, pos=0.4] {\iflangja{I/O・イベント完了}{I/O or event completion}} (rdy);
\end{tikzpicture}

  \caption{Process states and the transitions between them.}
  \label{fig:l02-states}
\end{figure}

\caution{A waiting process whose event completes does not resume running
immediately: it moves to the ready state and must be dispatched again.}
\Cref{tab:l02-trace} traces one process through its life cycle.  Note that
only one transition, the dispatch, gives the process the CPU; every other
transition takes it away or leaves the process without it.

\begin{table}[htbp]
  \centering
  \small
  \caption{State trace of a process that runs, is preempted once, reads a
    file, and exits.}
  \label{tab:l02-trace}
  \begin{tabular}{@{}llll@{}}
    \toprule
    Event & Transition & State after & Where it is \\
    \midrule
    process is being created     & ---                     & new        & --- \\
    PCB initialized, accepted    & admitted                & ready      & ready queue \\
    scheduler selects it         & scheduler dispatch      & running    & CPU \\
    its time on the CPU expires  & interrupt               & ready      & ready queue \\
    scheduler selects it         & scheduler dispatch      & running    & CPU \\
    it issues a disk read        & I/O or event wait       & waiting    & disk I/O queue \\
    the disk read completes      & I/O or event completion & ready      & ready queue \\
    scheduler selects it         & scheduler dispatch      & running    & CPU \\
    it calls exit                & exit                    & terminated & --- \\
    \bottomrule
  \end{tabular}
\end{table}

\section{Process creation}
\label{sec:l02-creation}

Processes are created by other processes.\source{P2L1, process life cycle:
creation; OSTEP ch.~5; Silberschatz et al.\ ch.~3.}  When the system boots
and the operating system has been loaded, it creates a few initial
processes, some of them with privileged access.  When a user logs in, a
shell process is created for the user, and every command typed into the
shell, for instance \texttt{ls} or \texttt{emacs}, runs as a new process
created by the shell.  A process that creates another process is its
\term{parent process}[親プロセス], and the process it creates is its
\term{child process}[子プロセス].  Since every process except the first
has a parent, the processes of a system form a tree rooted in a single
process (\cref{fig:l02-creation}\,(a)).

\begin{figure}[htbp]
  \centering
  % Left: a process tree.  Right: fork followed by exec.
% Usage: % Left: a process tree.  Right: fork followed by exec.
% Usage: % Left: a process tree.  Right: fork followed by exec.
% Usage: \input{figures/l02-creation}   (width ~116 mm)
\begin{tikzpicture}[x=1mm, y=1mm, font=\footnotesize\sffamily,
  pr/.style={draw, rounded corners=2pt, fill=ln-box, minimum width=13mm,
             minimum height=6mm, inner sep=1pt},
  bx/.style={draw, rounded corners=2pt, fill=white, text width=25mm, align=center,
             minimum height=9mm, inner sep=1.5pt, font=\scriptsize\sffamily},
  >={Stealth[length=1.6mm]},
  lab/.style={font=\scriptsize\sffamily, align=center},
  ttl/.style={font=\footnotesize\sffamily\itshape}]
  % ---- (a) process tree
  \node[pr, fill=ln-accent!22] (init) at (24,36) {\texttt{init}};
  \node[pr] (sh)  at (12,22) {\iflangja{シェル}{shell}};
  \node[pr] (d)   at (36,22) {\iflangja{デーモン}{daemon}};
  \node[pr] (ls)  at (7,8)   {\texttt{ls}};
  \node[pr] (em)  at (22,8)  {\texttt{emacs}};
  \draw[->] (init) -- (sh);
  \draw[->] (init) -- (d);
  \draw[->] (sh) -- (ls);
  \draw[->] (sh) -- (em);
  \node[ttl] at (21,-3) {\iflangja{(a) プロセスの木}{(a) process tree}};
  % ---- (b) fork, then exec in the child
  \node[bx] (par) at (66,34)
    {\iflangja{親：プログラム A\\PC = fork の直後}{parent: program A\\PC = after \texttt{fork}}};
  \node[bx] (ch)  at (66,9)
    {\iflangja{子：A の複製\\PC = fork の直後}{child: copy of A\\PC = after \texttt{fork}}};
  \node[bx] (ch2) at (103,9)
    {\iflangja{子：プログラム B\\PC = B の先頭命令}{child: program B\\PC = first instr.\ of B}};
  \draw[->, thick] (par) -- node[lab, right] {\texttt{fork}} (ch);
  \draw[->, thick] (ch) -- node[lab, above] {\texttt{exec}} (ch2);
  \draw[->, thick] (par.east) -- node[lab, above] {\iflangja{親は実行を続ける}{parent continues}} (117,34);
  \node[ttl] at (84,-3) {\iflangja{(b) 新しいプログラムの起動}{(b) starting a new program}};
\end{tikzpicture}
   (width ~116 mm)
\begin{tikzpicture}[x=1mm, y=1mm, font=\footnotesize\sffamily,
  pr/.style={draw, rounded corners=2pt, fill=ln-box, minimum width=13mm,
             minimum height=6mm, inner sep=1pt},
  bx/.style={draw, rounded corners=2pt, fill=white, text width=25mm, align=center,
             minimum height=9mm, inner sep=1.5pt, font=\scriptsize\sffamily},
  >={Stealth[length=1.6mm]},
  lab/.style={font=\scriptsize\sffamily, align=center},
  ttl/.style={font=\footnotesize\sffamily\itshape}]
  % ---- (a) process tree
  \node[pr, fill=ln-accent!22] (init) at (24,36) {\texttt{init}};
  \node[pr] (sh)  at (12,22) {\iflangja{シェル}{shell}};
  \node[pr] (d)   at (36,22) {\iflangja{デーモン}{daemon}};
  \node[pr] (ls)  at (7,8)   {\texttt{ls}};
  \node[pr] (em)  at (22,8)  {\texttt{emacs}};
  \draw[->] (init) -- (sh);
  \draw[->] (init) -- (d);
  \draw[->] (sh) -- (ls);
  \draw[->] (sh) -- (em);
  \node[ttl] at (21,-3) {\iflangja{(a) プロセスの木}{(a) process tree}};
  % ---- (b) fork, then exec in the child
  \node[bx] (par) at (66,34)
    {\iflangja{親：プログラム A\\PC = fork の直後}{parent: program A\\PC = after \texttt{fork}}};
  \node[bx] (ch)  at (66,9)
    {\iflangja{子：A の複製\\PC = fork の直後}{child: copy of A\\PC = after \texttt{fork}}};
  \node[bx] (ch2) at (103,9)
    {\iflangja{子：プログラム B\\PC = B の先頭命令}{child: program B\\PC = first instr.\ of B}};
  \draw[->, thick] (par) -- node[lab, right] {\texttt{fork}} (ch);
  \draw[->, thick] (ch) -- node[lab, above] {\texttt{exec}} (ch2);
  \draw[->, thick] (par.east) -- node[lab, above] {\iflangja{親は実行を続ける}{parent continues}} (117,34);
  \node[ttl] at (84,-3) {\iflangja{(b) 新しいプログラムの起動}{(b) starting a new program}};
\end{tikzpicture}
   (width ~116 mm)
\begin{tikzpicture}[x=1mm, y=1mm, font=\footnotesize\sffamily,
  pr/.style={draw, rounded corners=2pt, fill=ln-box, minimum width=13mm,
             minimum height=6mm, inner sep=1pt},
  bx/.style={draw, rounded corners=2pt, fill=white, text width=25mm, align=center,
             minimum height=9mm, inner sep=1.5pt, font=\scriptsize\sffamily},
  >={Stealth[length=1.6mm]},
  lab/.style={font=\scriptsize\sffamily, align=center},
  ttl/.style={font=\footnotesize\sffamily\itshape}]
  % ---- (a) process tree
  \node[pr, fill=ln-accent!22] (init) at (24,36) {\texttt{init}};
  \node[pr] (sh)  at (12,22) {\iflangja{シェル}{shell}};
  \node[pr] (d)   at (36,22) {\iflangja{デーモン}{daemon}};
  \node[pr] (ls)  at (7,8)   {\texttt{ls}};
  \node[pr] (em)  at (22,8)  {\texttt{emacs}};
  \draw[->] (init) -- (sh);
  \draw[->] (init) -- (d);
  \draw[->] (sh) -- (ls);
  \draw[->] (sh) -- (em);
  \node[ttl] at (21,-3) {\iflangja{(a) プロセスの木}{(a) process tree}};
  % ---- (b) fork, then exec in the child
  \node[bx] (par) at (66,34)
    {\iflangja{親：プログラム A\\PC = fork の直後}{parent: program A\\PC = after \texttt{fork}}};
  \node[bx] (ch)  at (66,9)
    {\iflangja{子：A の複製\\PC = fork の直後}{child: copy of A\\PC = after \texttt{fork}}};
  \node[bx] (ch2) at (103,9)
    {\iflangja{子：プログラム B\\PC = B の先頭命令}{child: program B\\PC = first instr.\ of B}};
  \draw[->, thick] (par) -- node[lab, right] {\texttt{fork}} (ch);
  \draw[->, thick] (ch) -- node[lab, above] {\texttt{exec}} (ch2);
  \draw[->, thick] (par.east) -- node[lab, above] {\iflangja{親は実行を続ける}{parent continues}} (117,34);
  \node[ttl] at (84,-3) {\iflangja{(b) 新しいプログラムの起動}{(b) starting a new program}};
\end{tikzpicture}

  \caption{(a) Processes form a tree rooted in the initial process.  (b) A
    child created by \texttt{fork} is a copy of its parent; \texttt{exec}
    replaces its program with a new one.}
  \label{fig:l02-creation}
\end{figure}

Most operating systems provide two mechanisms for creating processes, in
Unix called \texttt{fork} and \texttt{exec}:
\begin{description}
  \item[\texttt{fork}] With the \term[fork]{\texttt{fork}}[fork] system
    call the operating system creates a new PCB for the child, copies the
    PCB of the parent into it, and gives the child a copy of the parent's
    address space.  The child is a duplicate of the parent, with the same
    program and the same program counter, so after the call both parent and
    child continue at the instruction that follows \texttt{fork}.
  \item[\texttt{exec}] With the \term[exec]{\texttt{exec}}[exec] system
    call a process replaces its own image with a new program.  The
    operating system loads the new program into the address space of the
    process and reinitializes its PCB; the program counter now points to
    the first instruction of the new program.
\end{description}

To create a process that runs a \emph{different} program, a parent calls
\texttt{fork}, and the child then calls \texttt{exec}
(\cref{fig:l02-creation}\,(b)).  Since parent and child return from the
same \texttt{fork} to the same place, they tell each other apart by its
return value: \texttt{fork} returns 0 in the child and the child's process
identifier in the parent, or $-1$ in the parent if no child could be
created.

Every process tree has a root.  On Unix-based systems it is
\term[init]{\texttt{init}}[init], the first user process that the kernel
starts after booting, and every other process descends from it.  (On many
current Linux distributions the program that runs as this first process is
\texttt{systemd}.)  Android uses the same mechanism to start applications
quickly.  Its application processes descend from
\term{Zygote}[Zygote], a daemon started during boot whose sole purpose is to
launch applications; it loads the runtime and commonly used libraries once.
To launch an application, Zygote forks, and the child, which already
contains everything Zygote has loaded, becomes the process of the
application.

\needspace{18\baselineskip}
\begin{example}
  A program that creates a child process with \texttt{fork}, runs
  \texttt{ls -l} in it with \texttt{exec}, and waits until the child has
  finished:
  \begin{lstlisting}[language=C, numbers=none]
#include <sys/types.h>
#include <sys/wait.h>
#include <unistd.h>

int main(void) {
    pid_t pid = fork();
    if (pid == 0) {              /* child */
        execlp("ls", "ls", "-l", (char *)NULL);
        _exit(127);              /* only if exec failed */
    } else if (pid > 0) {        /* parent */
        waitpid(pid, NULL, 0);   /* wait for the child */
    }
    return 0;
}
  \end{lstlisting}
  After a successful \texttt{execlp} the child executes the code of
  \texttt{ls}; the call never returns.
\end{example}

\section{The CPU scheduler}
\label{sec:l02-scheduler}

For the CPU to start executing a process, the process must be ready, and at
any time several processes may be ready.\source{P2L1, role of the CPU
scheduler; OSTEP ch.~6--7.}  The operating system must decide which of them
runs next, and for how long.

\begin{definition}[CPU scheduler]
  The \term{CPU scheduler}[CPUスケジューラ] is the component of the
  operating system that manages how processes use the CPU.  It decides which
  of the ready processes is dispatched to the CPU and how long it may run.
\end{definition}

The processes that are ready wait for the CPU in the
\term{ready queue}[レディキュー].  To give the CPU to one of them, the
operating system performs three operations:
\begin{description}
  \item[Preempt] the running process is interrupted and its context is
    saved in its PCB; this is called \term{preemption}[プリエンプション];
  \item[Schedule] the scheduler runs its algorithm to choose the next
    process from the ready queue;
  \item[Dispatch] the operating system switches into the context of the
    chosen process and lets it run on the CPU; this is the
    \term{dispatch}[ディスパッチ] of the process.
\end{description}

The CPU is the resource that the processes need, so it should spend as much
of its time as possible running them rather than running the operating
system.  The scheduler therefore needs efficient algorithms, efficient
implementations of them, and efficient data structures to represent the
waiting processes and their scheduling information, the ready queue in
particular.

\subsection{How often to run the scheduler}

The more often the scheduler runs, the more CPU time it takes away from the
processes.\source{P2L1, length of process; OSTEP ch.~7.}  Suppose each
process runs for a time $T_p$ before the scheduler, which takes
$t_{\mathrm{sched}}$, chooses the next one (\cref{fig:l02-timeslice}).  Over
$n$ such rounds the fraction of the CPU time spent on useful work is
\begin{equation}
  \text{useful CPU work}
  \;=\; \frac{\text{total processing time}}{\text{total time}}
  \;=\; \frac{n\,T_p}{n\,T_p + n\,t_{\mathrm{sched}}}
  \;=\; \frac{T_p}{T_p + t_{\mathrm{sched}}} .
  \label{eq:l02-useful}
\end{equation}

\begin{figure}[htbp]
  \centering
  % Alternation of process timeslices and scheduler runs.
% Usage: % Alternation of process timeslices and scheduler runs.
% Usage: % Alternation of process timeslices and scheduler runs.
% Usage: \input{figures/l02-timeslice}   (width ~117 mm incl. the CPU label)
\begin{tikzpicture}[x=1mm, y=1mm, font=\footnotesize\sffamily,
  dim/.style={{Bar[width=2mm]}-{Bar[width=2mm]}, thin},
  pr/.style={draw, fill=ln-accent!22},
  sc/.style={draw, fill=ln-box}]
  \node[anchor=east] at (-1,3) {CPU};
  \draw[pr] (0,0)  rectangle (30,6)  node[midway] {P1};
  \draw[sc] (30,0) rectangle (38,6)  node[midway, font=\scriptsize\sffamily] {\iflangja{S}{S}};
  \draw[pr] (38,0) rectangle (68,6)  node[midway] {P2};
  \draw[sc] (68,0) rectangle (76,6)  node[midway, font=\scriptsize\sffamily] {\iflangja{S}{S}};
  \draw[pr] (76,0) rectangle (106,6) node[midway] {P3};
  \draw[sc] (106,0) rectangle (109,6);
  \foreach \a/\b in {0/30, 38/68, 76/106} {
    \draw[dim] (\a,-3) -- (\b,-3);
    \node[below] at ({(\a+\b)/2},-3.5) {$T_p$};
  }
  \foreach \a/\b in {30/38, 68/76} {
    \draw[dim] (\a,9) -- (\b,9);
    \node[above] at ({(\a+\b)/2},9.5) {$t_{\mathrm{sched}}$};
  }
\end{tikzpicture}
   (width ~117 mm incl. the CPU label)
\begin{tikzpicture}[x=1mm, y=1mm, font=\footnotesize\sffamily,
  dim/.style={{Bar[width=2mm]}-{Bar[width=2mm]}, thin},
  pr/.style={draw, fill=ln-accent!22},
  sc/.style={draw, fill=ln-box}]
  \node[anchor=east] at (-1,3) {CPU};
  \draw[pr] (0,0)  rectangle (30,6)  node[midway] {P1};
  \draw[sc] (30,0) rectangle (38,6)  node[midway, font=\scriptsize\sffamily] {\iflangja{S}{S}};
  \draw[pr] (38,0) rectangle (68,6)  node[midway] {P2};
  \draw[sc] (68,0) rectangle (76,6)  node[midway, font=\scriptsize\sffamily] {\iflangja{S}{S}};
  \draw[pr] (76,0) rectangle (106,6) node[midway] {P3};
  \draw[sc] (106,0) rectangle (109,6);
  \foreach \a/\b in {0/30, 38/68, 76/106} {
    \draw[dim] (\a,-3) -- (\b,-3);
    \node[below] at ({(\a+\b)/2},-3.5) {$T_p$};
  }
  \foreach \a/\b in {30/38, 68/76} {
    \draw[dim] (\a,9) -- (\b,9);
    \node[above] at ({(\a+\b)/2},9.5) {$t_{\mathrm{sched}}$};
  }
\end{tikzpicture}
   (width ~117 mm incl. the CPU label)
\begin{tikzpicture}[x=1mm, y=1mm, font=\footnotesize\sffamily,
  dim/.style={{Bar[width=2mm]}-{Bar[width=2mm]}, thin},
  pr/.style={draw, fill=ln-accent!22},
  sc/.style={draw, fill=ln-box}]
  \node[anchor=east] at (-1,3) {CPU};
  \draw[pr] (0,0)  rectangle (30,6)  node[midway] {P1};
  \draw[sc] (30,0) rectangle (38,6)  node[midway, font=\scriptsize\sffamily] {\iflangja{S}{S}};
  \draw[pr] (38,0) rectangle (68,6)  node[midway] {P2};
  \draw[sc] (68,0) rectangle (76,6)  node[midway, font=\scriptsize\sffamily] {\iflangja{S}{S}};
  \draw[pr] (76,0) rectangle (106,6) node[midway] {P3};
  \draw[sc] (106,0) rectangle (109,6);
  \foreach \a/\b in {0/30, 38/68, 76/106} {
    \draw[dim] (\a,-3) -- (\b,-3);
    \node[below] at ({(\a+\b)/2},-3.5) {$T_p$};
  }
  \foreach \a/\b in {30/38, 68/76} {
    \draw[dim] (\a,9) -- (\b,9);
    \node[above] at ({(\a+\b)/2},9.5) {$t_{\mathrm{sched}}$};
  }
\end{tikzpicture}

  \caption{Processes run for $T_p$ each; the scheduler (S) runs for
    $t_{\mathrm{sched}}$ between them.}
  \label{fig:l02-timeslice}
\end{figure}

The time $T_p$ that a process is allocated on the CPU is its
\term{timeslice}[タイムスライス].  The longer the timeslice relative to
$t_{\mathrm{sched}}$, the larger the useful fraction.  If the timeslice is
only as long as a run of the scheduler, half of the CPU time is lost; with a
timeslice ten times as long, the useful fraction is $10/11 \approx
\SI{91}{\percent}$.

\needspace{8\baselineskip}
\begin{example}
  Let $t_{\mathrm{sched}} = \SI{0.5}{\milli\second}$.  With a timeslice of
  \SI{2}{\milli\second}, \cref{eq:l02-useful} gives $2/2.5 =
  \SI{80}{\percent}$ useful work; with a timeslice of
  \SI{20}{\milli\second}, $20/20.5 \approx \SI{97.6}{\percent}$.  The longer
  timeslice has a price.  With ten ready processes, a process that has just
  used up its timeslice waits for the timeslices of the nine others and for
  ten runs of the scheduler, $9\times 20 + 10\times 0.5 =
  \SI{185}{\milli\second}$, before it runs again; with the
  \SI{2}{\milli\second} timeslice it would wait only
  $9\times 2 + 10\times 0.5 = \SI{23}{\milli\second}$.
\end{example}

Choosing a scheduler thus requires design decisions: what timeslice values
are appropriate, and which metrics decide the process to run next.  Lesson~7
examines these questions.

\subsection{Processes and I/O}

A running process may issue an I/O request, for example to read from
disk.\source{P2L1, what about I/O; Silberschatz et al.\ ch.~3.}  The
operating system delivers the request to the device and places the process
on the \emph{I/O queue} associated with that device.  The process remains
waiting there until the device completes the operation and responds to the
request.  The process is then ready again: depending on the load of the
system it is placed on the ready queue, or, if no other process is waiting,
it may be dispatched to the CPU right away.

\Cref{fig:l02-queues} shows all the ways in which a process that has left
the CPU finds its way back to the ready queue:
\begin{itemize}
  \item its I/O request has been served by the device;
  \item its timeslice has expired;
  \item it is a new process, created by a \texttt{fork};
  \item it was waiting for an interrupt, and the interrupt has occurred.
\end{itemize}

\begin{figure}[htbp]
  \centering
  % Queueing diagram: how processes move between the ready queue, the CPU
% and I/O queues.
% Usage: % Queueing diagram: how processes move between the ready queue, the CPU
% and I/O queues.
% Usage: % Queueing diagram: how processes move between the ready queue, the CPU
% and I/O queues.
% Usage: \input{figures/l02-queues}   (width ~108 mm)
\begin{tikzpicture}[x=1mm, y=1mm, font=\footnotesize\sffamily,
  >={Stealth[length=1.6mm]},
  res/.style={draw, circle, fill=ln-accent!22, minimum size=10mm, inner sep=0pt},
  bx/.style={draw, rounded corners=2pt, fill=white, minimum height=6mm,
             minimum width=22mm, inner sep=1pt, font=\scriptsize\sffamily},
  lab/.style={font=\scriptsize\sffamily}]
  % queue drawing: rectangle with slots, open at the left
  \newcommand{\lnqueue}[3]{% x, y, label
    \draw[fill=ln-box] (#1,#2-3) rectangle ++(28,6);
    \foreach \i in {1,...,4} \draw (#1+28-4*\i,#2-3) -- ++(0,6);
    \node[lab, anchor=south] at (#1+14,#2+3) {#3};}
  \lnqueue{16}{40}{\iflangja{レディキュー}{ready queue}}
  \node[res] (cpu) at (64,40) {CPU};
  \draw[->] (44,40) -- (cpu.west);
  % right-hand bus and left-hand return line
  \draw (cpu.east) -- (102,40) -- (102,-8);
  \draw[->] (2,-8) -- (2,40) -- (16,40);
  % I/O
  \draw[->] (102,26) -- node[lab, above] {\iflangja{I/O 要求}{I/O request}} (80,26);
  \lnqueue{52}{26}{\iflangja{I/O キュー}{I/O queue}}
  \node[res, minimum size=8mm] (io) at (36,26) {I/O};
  \draw[->] (52,26) -- (io.east);
  \draw (io.west) -- (2,26);
  % timeslice expired
  \draw (102,14) -- node[lab, above, pos=0.5] {\iflangja{タイムスライス満了}{timeslice expired}} (2,14);
  % fork
  \draw[->] (102,3) -- node[lab, above] {\iflangja{子を fork}{fork a child}} (74,3);
  \node[bx] (ch) at (63,3) {\iflangja{子プロセスが作られる}{child is created}};
  \draw (ch.west) -- (2,3);
  % interrupt
  \draw[->] (102,-8) -- node[lab, above] {\iflangja{割り込みを待つ}{wait for an interrupt}} (74,-8);
  \node[bx] (in) at (63,-8) {\iflangja{割り込みが発生}{interrupt occurs}};
  \draw (in.west) -- (2,-8);
\end{tikzpicture}
   (width ~108 mm)
\begin{tikzpicture}[x=1mm, y=1mm, font=\footnotesize\sffamily,
  >={Stealth[length=1.6mm]},
  res/.style={draw, circle, fill=ln-accent!22, minimum size=10mm, inner sep=0pt},
  bx/.style={draw, rounded corners=2pt, fill=white, minimum height=6mm,
             minimum width=22mm, inner sep=1pt, font=\scriptsize\sffamily},
  lab/.style={font=\scriptsize\sffamily}]
  % queue drawing: rectangle with slots, open at the left
  \newcommand{\lnqueue}[3]{% x, y, label
    \draw[fill=ln-box] (#1,#2-3) rectangle ++(28,6);
    \foreach \i in {1,...,4} \draw (#1+28-4*\i,#2-3) -- ++(0,6);
    \node[lab, anchor=south] at (#1+14,#2+3) {#3};}
  \lnqueue{16}{40}{\iflangja{レディキュー}{ready queue}}
  \node[res] (cpu) at (64,40) {CPU};
  \draw[->] (44,40) -- (cpu.west);
  % right-hand bus and left-hand return line
  \draw (cpu.east) -- (102,40) -- (102,-8);
  \draw[->] (2,-8) -- (2,40) -- (16,40);
  % I/O
  \draw[->] (102,26) -- node[lab, above] {\iflangja{I/O 要求}{I/O request}} (80,26);
  \lnqueue{52}{26}{\iflangja{I/O キュー}{I/O queue}}
  \node[res, minimum size=8mm] (io) at (36,26) {I/O};
  \draw[->] (52,26) -- (io.east);
  \draw (io.west) -- (2,26);
  % timeslice expired
  \draw (102,14) -- node[lab, above, pos=0.5] {\iflangja{タイムスライス満了}{timeslice expired}} (2,14);
  % fork
  \draw[->] (102,3) -- node[lab, above] {\iflangja{子を fork}{fork a child}} (74,3);
  \node[bx] (ch) at (63,3) {\iflangja{子プロセスが作られる}{child is created}};
  \draw (ch.west) -- (2,3);
  % interrupt
  \draw[->] (102,-8) -- node[lab, above] {\iflangja{割り込みを待つ}{wait for an interrupt}} (74,-8);
  \node[bx] (in) at (63,-8) {\iflangja{割り込みが発生}{interrupt occurs}};
  \draw (in.west) -- (2,-8);
\end{tikzpicture}
   (width ~108 mm)
\begin{tikzpicture}[x=1mm, y=1mm, font=\footnotesize\sffamily,
  >={Stealth[length=1.6mm]},
  res/.style={draw, circle, fill=ln-accent!22, minimum size=10mm, inner sep=0pt},
  bx/.style={draw, rounded corners=2pt, fill=white, minimum height=6mm,
             minimum width=22mm, inner sep=1pt, font=\scriptsize\sffamily},
  lab/.style={font=\scriptsize\sffamily}]
  % queue drawing: rectangle with slots, open at the left
  \newcommand{\lnqueue}[3]{% x, y, label
    \draw[fill=ln-box] (#1,#2-3) rectangle ++(28,6);
    \foreach \i in {1,...,4} \draw (#1+28-4*\i,#2-3) -- ++(0,6);
    \node[lab, anchor=south] at (#1+14,#2+3) {#3};}
  \lnqueue{16}{40}{\iflangja{レディキュー}{ready queue}}
  \node[res] (cpu) at (64,40) {CPU};
  \draw[->] (44,40) -- (cpu.west);
  % right-hand bus and left-hand return line
  \draw (cpu.east) -- (102,40) -- (102,-8);
  \draw[->] (2,-8) -- (2,40) -- (16,40);
  % I/O
  \draw[->] (102,26) -- node[lab, above] {\iflangja{I/O 要求}{I/O request}} (80,26);
  \lnqueue{52}{26}{\iflangja{I/O キュー}{I/O queue}}
  \node[res, minimum size=8mm] (io) at (36,26) {I/O};
  \draw[->] (52,26) -- (io.east);
  \draw (io.west) -- (2,26);
  % timeslice expired
  \draw (102,14) -- node[lab, above, pos=0.5] {\iflangja{タイムスライス満了}{timeslice expired}} (2,14);
  % fork
  \draw[->] (102,3) -- node[lab, above] {\iflangja{子を fork}{fork a child}} (74,3);
  \node[bx] (ch) at (63,3) {\iflangja{子プロセスが作られる}{child is created}};
  \draw (ch.west) -- (2,3);
  % interrupt
  \draw[->] (102,-8) -- node[lab, above] {\iflangja{割り込みを待つ}{wait for an interrupt}} (74,-8);
  \node[bx] (in) at (63,-8) {\iflangja{割り込みが発生}{interrupt occurs}};
  \draw (in.west) -- (2,-8);
\end{tikzpicture}

  \caption{Paths of a process between the ready queue, the CPU and the
    places where it waits.  Every path returns to the ready queue.}
  \label{fig:l02-queues}
\end{figure}

The responsibilities are divided accordingly.  The scheduler manages the
ready queue, and it decides when to switch contexts and to which process.
It does not manage the I/O queues, which belong to the handling of the
devices, and it does not produce the events that waiting processes wait for:
those come from devices, timers or other processes.  The scheduler takes
over only once a process is on the ready queue again.

\begin{keypoint}
  The scheduler preempts, schedules and dispatches processes that are in the
  ready queue.  Because it runs on the same CPU, its own time is overhead
  (\cref{eq:l02-useful}); the timeslice trades this overhead against how long
  ready processes wait.
\end{keypoint}

\section{Inter-process communication}
\label{sec:l02-ipc}

Many applications are structured as several processes that must interact.
A web application, for example, may consist of a web server process that
handles the requests of clients and a database process that stores the
data.\source{P2L1, can processes interact; Silberschatz et al.\ ch.~3.}
The operating system, however, isolates processes from one another (Lesson~1):
each has its own address space, and no process can access the memory of
another.

\begin{definition}[Inter-process communication]
  The mechanisms that the operating system provides for processes to
  interact are called
  \term{inter-process communication}[プロセス間通信]\index{IPC|see{inter-process communication}}
  (IPC).  IPC mechanisms transfer data and information
  between address spaces, maintain the protection and isolation of the
  processes, and give applications flexibility and good performance.
\end{definition}

There are two basic kinds of IPC (\cref{fig:l02-ipc}):
\begin{description}
  \item[Message passing] In \term{message passing}[メッセージパッシング]
    the operating system establishes a communication channel, such as a
    shared buffer in kernel memory.  Processes send messages by writing into
    the channel and receive messages by reading from it, using system calls.
    The operating system manages the channel and offers the same
    well-defined interface, such as \texttt{send} and \texttt{recv}, to all
    applications.  The price is overhead: every piece of information is
    copied from the address space of the sender into the channel in the
    kernel, and from there into the address space of the receiver, with a
    crossing into the kernel on each side.
  \item[Shared memory] With \term{shared memory}[共有メモリ] the operating
    system establishes a region of physical memory and maps it into the
    address spaces of both processes.  The processes then read and write
    this memory directly, as they would their own, and the operating system
    is out of the way.  The price is that the operating system no longer
    defines how the memory is used.  There is no fixed, well-defined
    interface: the processes must agree on the format of the data and
    coordinate their accesses themselves, developers may have to
    re-implement this code for every application, and it is more
    error-prone.
\end{description}

\begin{figure}[htbp]
  \centering
  % Message-passing IPC versus shared-memory IPC.
% Usage: % Message-passing IPC versus shared-memory IPC.
% Usage: % Message-passing IPC versus shared-memory IPC.
% Usage: \input{figures/l02-ipc}   (width ~116 mm)
\begin{tikzpicture}[x=1mm, y=1mm, font=\footnotesize\sffamily,
  >={Stealth[length=1.6mm]},
  lab/.style={font=\scriptsize\sffamily, align=center},
  pname/.style={font=\footnotesize\sffamily\bfseries, anchor=north},
  ttl/.style={font=\footnotesize\sffamily\itshape}]
  % ---- (a) message passing
  \draw[rounded corners=2pt, fill=white] (0,23) rectangle (20,38);
  \node[pname] at (10,38) {P1};
  \draw[rounded corners=2pt, fill=white] (34,23) rectangle (54,38);
  \node[pname] at (44,38) {P2};
  \draw[fill=ln-box] (0,0) rectangle (54,13);
  \node[lab, anchor=south west, text=ln-muted] at (0.5,0.3) {\iflangja{カーネル}{kernel}};
  \draw[fill=ln-accent!22] (16,4) rectangle (38,10) node[midway, lab] {\iflangja{チャネル}{channel}};
  \draw[->, thick] (14,23) -- (20,10);
  \draw[->, thick] (34,10) -- (40,23);
  \node[lab, anchor=east] at (13.5,17.5) {\texttt{send}\\\iflangja{（コピー）}{(copy)}};
  \node[lab, anchor=west] at (40.5,17.5) {\texttt{recv}\\\iflangja{（コピー）}{(copy)}};
  \node[ttl] at (27,-5) {\iflangja{(a) メッセージパッシング}{(a) message passing}};
  % ---- (b) shared memory
  \draw[rounded corners=2pt, fill=white] (62,23) rectangle (82,38);
  \node[pname] at (72,38) {P1};
  \draw[rounded corners=2pt, fill=white] (96,23) rectangle (116,38);
  \node[pname] at (106,38) {P2};
  \draw[fill=ln-accent!22] (66,25) rectangle (78,30);
  \draw[fill=ln-accent!22] (100,28) rectangle (112,33);
  \draw[fill=ln-box] (62,0) rectangle (116,13);
  \node[lab, anchor=south west, text=ln-muted] at (62.5,0.3) {\iflangja{物理メモリ}{physical memory}};
  \draw[fill=ln-accent!22] (80,4) rectangle (98,10) node[midway, lab] {\iflangja{共有領域}{shared}};
  \draw[thick, densely dashed] (72,25) -- (85,10);
  \draw[thick, densely dashed] (106,28) -- (93,10);
  \node[lab, text=ln-muted] at (89,19) {\iflangja{対応付け}{mapped}};
  \node[ttl] at (89,-5) {\iflangja{(b) 共有メモリ}{(b) shared memory}};
\end{tikzpicture}
   (width ~116 mm)
\begin{tikzpicture}[x=1mm, y=1mm, font=\footnotesize\sffamily,
  >={Stealth[length=1.6mm]},
  lab/.style={font=\scriptsize\sffamily, align=center},
  pname/.style={font=\footnotesize\sffamily\bfseries, anchor=north},
  ttl/.style={font=\footnotesize\sffamily\itshape}]
  % ---- (a) message passing
  \draw[rounded corners=2pt, fill=white] (0,23) rectangle (20,38);
  \node[pname] at (10,38) {P1};
  \draw[rounded corners=2pt, fill=white] (34,23) rectangle (54,38);
  \node[pname] at (44,38) {P2};
  \draw[fill=ln-box] (0,0) rectangle (54,13);
  \node[lab, anchor=south west, text=ln-muted] at (0.5,0.3) {\iflangja{カーネル}{kernel}};
  \draw[fill=ln-accent!22] (16,4) rectangle (38,10) node[midway, lab] {\iflangja{チャネル}{channel}};
  \draw[->, thick] (14,23) -- (20,10);
  \draw[->, thick] (34,10) -- (40,23);
  \node[lab, anchor=east] at (13.5,17.5) {\texttt{send}\\\iflangja{（コピー）}{(copy)}};
  \node[lab, anchor=west] at (40.5,17.5) {\texttt{recv}\\\iflangja{（コピー）}{(copy)}};
  \node[ttl] at (27,-5) {\iflangja{(a) メッセージパッシング}{(a) message passing}};
  % ---- (b) shared memory
  \draw[rounded corners=2pt, fill=white] (62,23) rectangle (82,38);
  \node[pname] at (72,38) {P1};
  \draw[rounded corners=2pt, fill=white] (96,23) rectangle (116,38);
  \node[pname] at (106,38) {P2};
  \draw[fill=ln-accent!22] (66,25) rectangle (78,30);
  \draw[fill=ln-accent!22] (100,28) rectangle (112,33);
  \draw[fill=ln-box] (62,0) rectangle (116,13);
  \node[lab, anchor=south west, text=ln-muted] at (62.5,0.3) {\iflangja{物理メモリ}{physical memory}};
  \draw[fill=ln-accent!22] (80,4) rectangle (98,10) node[midway, lab] {\iflangja{共有領域}{shared}};
  \draw[thick, densely dashed] (72,25) -- (85,10);
  \draw[thick, densely dashed] (106,28) -- (93,10);
  \node[lab, text=ln-muted] at (89,19) {\iflangja{対応付け}{mapped}};
  \node[ttl] at (89,-5) {\iflangja{(b) 共有メモリ}{(b) shared memory}};
\end{tikzpicture}
   (width ~116 mm)
\begin{tikzpicture}[x=1mm, y=1mm, font=\footnotesize\sffamily,
  >={Stealth[length=1.6mm]},
  lab/.style={font=\scriptsize\sffamily, align=center},
  pname/.style={font=\footnotesize\sffamily\bfseries, anchor=north},
  ttl/.style={font=\footnotesize\sffamily\itshape}]
  % ---- (a) message passing
  \draw[rounded corners=2pt, fill=white] (0,23) rectangle (20,38);
  \node[pname] at (10,38) {P1};
  \draw[rounded corners=2pt, fill=white] (34,23) rectangle (54,38);
  \node[pname] at (44,38) {P2};
  \draw[fill=ln-box] (0,0) rectangle (54,13);
  \node[lab, anchor=south west, text=ln-muted] at (0.5,0.3) {\iflangja{カーネル}{kernel}};
  \draw[fill=ln-accent!22] (16,4) rectangle (38,10) node[midway, lab] {\iflangja{チャネル}{channel}};
  \draw[->, thick] (14,23) -- (20,10);
  \draw[->, thick] (34,10) -- (40,23);
  \node[lab, anchor=east] at (13.5,17.5) {\texttt{send}\\\iflangja{（コピー）}{(copy)}};
  \node[lab, anchor=west] at (40.5,17.5) {\texttt{recv}\\\iflangja{（コピー）}{(copy)}};
  \node[ttl] at (27,-5) {\iflangja{(a) メッセージパッシング}{(a) message passing}};
  % ---- (b) shared memory
  \draw[rounded corners=2pt, fill=white] (62,23) rectangle (82,38);
  \node[pname] at (72,38) {P1};
  \draw[rounded corners=2pt, fill=white] (96,23) rectangle (116,38);
  \node[pname] at (106,38) {P2};
  \draw[fill=ln-accent!22] (66,25) rectangle (78,30);
  \draw[fill=ln-accent!22] (100,28) rectangle (112,33);
  \draw[fill=ln-box] (62,0) rectangle (116,13);
  \node[lab, anchor=south west, text=ln-muted] at (62.5,0.3) {\iflangja{物理メモリ}{physical memory}};
  \draw[fill=ln-accent!22] (80,4) rectangle (98,10) node[midway, lab] {\iflangja{共有領域}{shared}};
  \draw[thick, densely dashed] (72,25) -- (85,10);
  \draw[thick, densely dashed] (106,28) -- (93,10);
  \node[lab, text=ln-muted] at (89,19) {\iflangja{対応付け}{mapped}};
  \node[ttl] at (89,-5) {\iflangja{(b) 共有メモリ}{(b) shared memory}};
\end{tikzpicture}

  \caption{(a) Message passing copies every message into a channel in the
    kernel and out again.  (b) With shared memory, the same physical pages
    are mapped into both address spaces.}
  \label{fig:l02-ipc}
\end{figure}

Whether shared memory performs better than message passing depends on the
communication pattern.  An individual exchange through shared memory is
cheap, because nothing is copied through the kernel.  Establishing the
shared mapping, however, is itself an expensive operation, and it pays off
only when its cost is amortized over a sufficiently large number of
exchanges.

\needspace{7\baselineskip}
\begin{example}
  Suppose that sending a message through a kernel channel costs
  \SI{6}{\micro\second} (two system calls and two copies), and that
  exchanging the same data through shared memory costs
  \SI{1}{\micro\second} after the mapping has been established at a
  one-time cost of \SI{300}{\micro\second}.  For $n$ messages, shared memory
  costs $300 + n$ and message passing $6n$ microseconds.  Shared memory is
  cheaper only when $300 + n < 6n$, that is, for more than 60 messages.
\end{example}

The IPC mechanisms of Unix-based systems and their use are the subject of
Lesson~9.

\begin{keypoint}[Lesson summary]
  A process is an executing program, represented by its address
  space---text, data, heap and stack, translated from virtual to physical
  addresses by a page table---and by the execution state that the operating
  system keeps in its PCB.  A context switch saves one process's state into
  its PCB and restores another's, at a direct cost and at the indirect cost
  of a cold cache.  Processes move between the new, ready, running, waiting
  and terminated states; they are created with \texttt{fork} and
  \texttt{exec}; the CPU scheduler picks the next ready process, with a
  timeslice that trades overhead against waiting time; and IPC, by message
  passing or shared memory, lets isolated processes communicate.
\end{keypoint}

\end{document}
